  inverting the Miura transform; the computation is
  triangular and algorithmic at each spin.
\end{enumerate}
\end{remark}

\begin{remark}[Effective central charge and intertwining
in the Miura basis]
\label{rem:spin2-ceff-miura}
The coproduct~\eqref{eq:coprod-T} in the
$\cW_{1+\infty}$ field basis has cross-term coefficient
$(\Psi - 1)/\Psi$ (not $1/\Psi$), as derived in step~4.
This coefficient controls two structural invariants of the
image $\Delta_z(T)$ in $V \otimes V$:
\begin{enumerate}[label=\textup{(\arabic*)}]
\item \emph{Effective central charge.}
  The image $\Delta_z(T_n)$ generates a Virasoro
  subalgebra in the tensor product mode algebra with
  effective central charge
  \begin{equation}\label{eq:c-eff-spin2}
    c_{\mathrm{eff}}(\Psi) = 2 + 2(\Psi - 1)^2.
  \end{equation}
  At $\Psi = 1$ (free boson): $c_{\mathrm{eff}} = 2$
  (two decoupled copies, the cross-term vanishes).
  At $\Psi = 2$: $c_{\mathrm{eff}} = 4$.
  The quadratic growth $c_{\mathrm{eff}} \sim 2\Psi^2$ at
  large $\Psi$ reflects the quadratic Miura nonlinearity
  $\psi_2 = T + J^2/(2\Psi)$.
  The formula is independent of the spectral parameter~$z$
  (the $z$-dependent terms in $\Delta_z(T)$ do not
  contribute to the vacuum expectation value of the
  commutator $[\Delta_z(T_2), \Delta_z(T_{-2})]$).
\item \emph{Intertwining with $\Delta_z(J)$.}
  The commutator
  $[\Delta_z(T_n),\, \Delta_z(J_m)]
  = -\Psi \cdot m \cdot \Delta_z(J_{n+m})$
  has intertwining factor $\Psi$ (not~$1$).
  The extra factor $\Psi - 1$ beyond the original algebra
  relation $[T_n, J_m] = -m\,J_{n+m}$ arises from the
  cross-term $(\Psi - 1)/\Psi \cdot (J \otimes J)_n$
  commuting with $\Delta_z(J_m) = J_m^L + J_m^R$:
  the left and right Heisenberg commutators each contribute
  $-((\Psi-1)/\Psi) \cdot \Psi \cdot m = -(\Psi - 1)m$.
  At $\Psi = 1$: factor~$1$ (algebra homomorphism).
\end{enumerate}
Both features are consequences of the Miura nonlinearity:
the Drinfeld coproduct is an algebra homomorphism of the
Yangian $Y(\widehat{\mathfrak{gl}}_1)$ (preserving the
$\psi$-generator relations), but when expressed in the
$\cW_{1+\infty}$ field basis via the nonlinear Miura
inversion $T = \psi_2 - J^2/(2\Psi)$, the image acquires
$\Psi$-dependent structural invariants that are absent in
the $\psi$-basis.
\end{remark}

\begin{remark}[Descent to $\cW_N$]
\label{rem:w-infty-descent}
Drinfeld--Sokolov reduction gives
$\pi_N \colon \cW_{1+\infty}[\Psi]
\twoheadrightarrow \cW_N$
for each $N \geq 2$. Since the spectral coproduct of
Theorem~\ref{thm:w-infty-chiral-qg} is defined
on the full $\cW_{1+\infty}$, it descends to $\cW_N$
provided $\pi_N \otimes \pi_N$ intertwines $\Delta_z$ with
a coproduct on $\cW_N$. For $\cW_{1+\infty}$ itself,
OPE compatibility holds at all spins by the $GL_1$
triviality argument (step~6 of the proof); for finite
$N$, the intertwining $(\pi_N \otimes \pi_N) \circ
\Delta_z = \Delta_z^{(N)} \circ \pi_N$ follows from the
compatibility of the Miura transform with truncation
(Remark~\ref{rem:w-infty-vertex-gap}).
The tower
$\cdots \twoheadrightarrow \cW_{N+1}
\twoheadrightarrow \cW_N
\twoheadrightarrow \cdots
\twoheadrightarrow \cW_2 = \mathrm{Vir}_c$
is a projective system of algebras with spectral
coproducts.
At $N = 2$: class $M$, $\kappa = c/2$, irregular KZ
connection. At $N = 3$: $\kappa(\cW_3) = 5c/6$ (where
$H_3 - 1 = 5/6$).
The Gaiotto--Rap\v{c}\'ak $Y$-algebras
$Y_{N_1,N_2,N_3}[\Psi]$~\cite{GR17} are truncations of
$\cW_{1+\infty}$; the $S_3$-triality is a symmetry
of the spectral coproduct, corresponding to three Hecke
correspondences on spiked-instanton moduli
(Rap\v{c}\'ak--Soibelman--Yang--Zhao~\cite{RSYZ20}).
The $\cD$-module $\cF_n^{\mathrm{ord}}(\cW_{1+\infty})$
has irregular singularities of unbounded Poincar\'e rank
(spin-$s$ contributes pole order $2s - 1$ after $d$-log
absorption), providing the boundary algebra of higher-spin
holographic gravity.
\end{remark}

\begin{remark}[Proof structure and remaining questions]
\label{rem:w-infty-vertex-gap}
Three comments on the proof of
Theorem~\textup{\ref{thm:w-infty-chiral-qg}}.

\smallskip
\noindent\emph{(1) Why RTT alone does not suffice.}
The RTT relation~\eqref{eq:rtt-gl1} is a constraint on the
\emph{associative algebra} structure (mode commutators).
The vertex bialgebra axiom~\eqref{eq:ope-compat} is a
constraint on the \emph{vertex algebra} structure (all
$n$-products $a_{(n)}b$, including the regular products for
$n < 0$). For $\mathfrak{gl}_1$ the RTT relation is
trivially satisfied because the $R$-matrix is scalar; this
confirms associative-algebra compatibility but does not, by
itself, establish the full vertex-algebra compatibility.
The two independent arguments of step~6 (the
coderivation property on the Koszul-locus bar complex, and
the JKL vertex bialgebra theorem on the CoHA) close this
gap.

\smallskip
\noindent\emph{(2) Truncation to $\cW_N$.}
Theorem~\ref{thm:glN-chiral-qg} resolves this question:
for the truncations $\cW_N$ at finite $N \geq 2$, the
$R$-matrix of $Y(\widehat{\mathfrak{gl}}_N)$ is Yang's
rational $R$-matrix $R(u) = u\, I + \Psi\, P$ (no longer
scalar), and the OPE compatibility at each truncation level
is established by the same two arguments.
The coderivation argument (argument~A) applies to
$\cW_N$ at generic level by the same Koszul-locus reasoning.
The JKL argument (argument~B) applies to
$Y(\widehat{\mathfrak{gl}}_N)$ by their Theorem~C for
general quivers.

\smallskip
\noindent\emph{(3) Antipode and Drinfeld double.}
The Schiffmann--Mellit--Minets--Vasserot
isomorphism~\cite{SMMV23} gives the positive half
$\cH_{\mathrm{Jor}} \cong U(\cW_{1+\infty}[\Psi])^+$.
The Drinfeld double
$D(\cH_{\mathrm{Jor}}) \cong
Y(\widehat{\mathfrak{gl}}_1)$ (Yang--Zhao~\cite{YZ18a})
provides the full algebra with antipode
$S(T(u)) = T(u)^{-1}$. Whether this lifts to a
vertex-algebraic antipode on $\cW_{1+\infty}$, yielding a
vertex Hopf algebra in the sense of Etingof--Kazhdan, is a
natural question beyond the scope of the present theorem.
\end{remark}

\begin{lemma}[Miura inversion of the spectral coproduct at
spin~$2$]
\label{lem:coprod-T-miura}
Let $\psi_1 = J$ and $\psi_2 = T + J^2/(2\Psi)$ be the
Miura identification of
Proch\'azka--Rap\v{c}\'ak~\textup{\cite{PR19}}.
The spectral coproduct
$\Delta_z(T)
= \Delta_z(\psi_2) - (1/(2\Psi))\,\Delta_z(J)^2$
equals
\begin{equation}\label{eq:coprod-T-derived}
  \Delta_z(T)
  \;=\;
  T \otimes 1
  + 1 \otimes T
  + \tfrac{\Psi - 1}{\Psi}\,
  J \otimes J
  + z\,(1 \otimes J).
\end{equation}
\end{lemma}

\begin{proof}
The Drinfeld formula
$\Delta_z(T(u)) = T(u) \otimes T(u-z)$
gives, at the $u^{-2}$ coefficient,
\[
  \Delta_z(\psi_2)
  \;=\;
  \psi_2 \otimes 1
  + 1 \otimes \psi_2
  + \psi_1 \otimes \psi_1
  + z\,(1 \otimes \psi_1).
\]
Since $\Delta_z$ is an algebra homomorphism and
$J = \psi_1$ is primitive,
\begin{align*}
  \Delta_z(J)^2
  &= (J \otimes 1 + 1 \otimes J)^2 \\
  &= J^2 \otimes 1
  + 2\,J \otimes J
  + 1 \otimes J^2.
\end{align*}
Substituting $T = \psi_2 - J^2/(2\Psi)$:
\begin{align*}
  \Delta_z(T)
  &= \Delta_z(\psi_2)
  - \tfrac{1}{2\Psi}\,\Delta_z(J)^2 \\
  &= \Bigl[
    \bigl(T + \tfrac{J^2}{2\Psi}\bigr) \otimes 1
    + 1 \otimes \bigl(T + \tfrac{J^2}{2\Psi}\bigr)
    + J \otimes J
    + z\,(1 \otimes J)
  \Bigr] \\
  &\qquad
  - \tfrac{1}{2\Psi}\Bigl[
    J^2 \otimes 1
    + 2\,J \otimes J
    + 1 \otimes J^2
  \Bigr] \\
  &= T \otimes 1
  + 1 \otimes T
  + \bigl(1 - \tfrac{1}{\Psi}\bigr)\,
  J \otimes J
  + z\,(1 \otimes J).
\end{align*}
Since $1 - 1/\Psi = (\Psi - 1)/\Psi$, this
is~\eqref{eq:coprod-T-derived}. The
$J^2/(2\Psi)$ terms cancel between the Miura
substitution and the $\Delta_z(J)^2$ correction.
At $\Psi = 1$ (free boson, $c = 1$) the
$J \otimes J$ cross-term vanishes and
$\Delta_z(T) = T \otimes 1 + 1 \otimes T
+ z\,(1 \otimes J)$;
in the classical limit $\Psi \to \infty$ the
coefficient tends to~$1$.
\end{proof}


% ----------------------------------------------------------------
\subsection{The chiral quantum group structure on
$\cW_N$ for $N \geq 2$}
\label{subsec:glN-chiral-qg}

Theorem~\ref{thm:w-infty-chiral-qg} established the chiral
quantum group structure on $\cW_{1+\infty}[\Psi]$ via the
scalar Maulik--Okounkov $R$-matrix. For the truncations
$\cW_N = \cW_{1+\infty}/\cI_N$ (generators of spins
$1, 2, \ldots, N$) the $R$-matrix is no longer scalar: it
is Yang's rational $R$-matrix on $\CC^N \otimes \CC^N$,
and the RTT relation is a non-trivial matrix identity.

\begin{theorem}[$Y(\widehat{\mathfrak{gl}}_N)$ chiral
quantum group for $N \geq 2$]
\label{thm:glN-chiral-qg}
For each $N \geq 2$ and generic $\Psi$, the truncation
$\cW_N[\Psi]$ carries a chiral quantum group datum in
the sense of
Theorem~\textup{\ref{thm:chiral-qg-equiv}}: an $N \times N$
matrix transfer matrix, Yang's rational $R$-matrix with
level prefix~$\Psi$, and a chiral coproduct with all
OPE compatibility verified.
The precise content is as follows.
\begin{enumerate}[label=\textup{(\Roman*)}]
\item \textup{($N \times N$ transfer matrix.)}
  The affine Yangian $Y(\widehat{\mathfrak{gl}}_N)$ has
  transfer matrix
  \begin{equation}\label{eq:glN-transfer}
    T(u)
    = I_N + \sum_{r \geq 1} T^{(r)}\, u^{-r},
    \qquad
    (T^{(r)})_{ij} = t_{ij}^{(r)},
    \quad 1 \leq i,j \leq N,
  \end{equation}
  an $N \times N$ matrix of formal power series in
  $u^{-1}$. The Drinfeld--Sokolov reduction
  $\pi_N \colon \cW_{1+\infty}[\Psi]
  \twoheadrightarrow \cW_N$ sends the scalar transfer
  matrix $T(u) = 1 + \sum \psi_n\, u^{-n}$ of
  Theorem~\textup{\ref{thm:w-infty-chiral-qg}}
  to the quantum Miura factorisation
  \begin{equation}\label{eq:glN-miura}
    T(u)
    = (u + \Lambda_1^{(1)})
    (u + \Lambda_2^{(1)})
    \cdots
    (u + \Lambda_N^{(1)}),
  \end{equation}
  where the $\Lambda_i^{(1)}$ are the Miura fields of
  Proch\'azka--Rap\v{c}\'ak~\textup{\cite{PR19}}.
\item \textup{(Yang's rational $R$-matrix.)}
  The $R$-matrix for $Y(\widehat{\mathfrak{gl}}_N)$
  in the fundamental representation is
  \begin{equation}\label{eq:glN-yang-r}
    % AP126: level prefix Psi mandatory; Psi=0 -> R = uI (trivial). Verified.
    R(u)
    = u\, I_{N^2} + \Psi\, P
    \quad \in \End(\CC^N \otimes \CC^N),
  \end{equation}
  where $P(e_i \otimes e_j) = e_j \otimes e_i$ is the
  permutation and $\Psi$ is the level.
  The classical $r$-matrix
  \textup{(}trace-form, level prefix
  $\Psi$\textup{)} is
  % AP126: Psi=0 -> r=0. Verified.
  \begin{equation}\label{eq:glN-classical-r}
    r(z) = \frac{\Psi\, P}{z},
  \end{equation}
  with $r|_{\Psi = 0} = 0$. The unitarity relation is
  $R(u)\, R(-u) = (\Psi^2 - u^2)\, I$.
  The Yang--Baxter equation
  \begin{equation}\label{eq:glN-ybe}
    R_{12}(u-v)\, R_{13}(u)\, R_{23}(v)
    = R_{23}(v)\, R_{13}(u)\, R_{12}(u-v)
  \end{equation}
  holds for all $N$ and all $\Psi$
  \textup{(}computed numerically at $N = 2, 3$ by
  \texttt{glN\_affine\_yangian\_chiral\_qg\_engine.py};
  proved algebraically in
  Molev~\textup{\cite{Molev07}},
  Chapter~$1$\textup{)}.
  At $N = 1$: $P = 1$ and $R(u) = u + \Psi$ is scalar,
  recovering the class-$G$ case of
  Theorem~\textup{\ref{thm:w-infty-chiral-qg}}.
\item \textup{(Drinfeld coproduct: matrix multiplication.)}
  The Drinfeld coproduct on $T(u)$ is
  \begin{equation}\label{eq:glN-drinfeld-coprod}
    \Delta_z(T(u)) = T(u) \cdot T(u - z)
    \qquad
    \textup{(matrix multiplication in $\CC^N$)},
  \end{equation}
  that is, in components:
  \begin{equation}\label{eq:glN-coprod-components}
    \Delta_z(t_{ij}(u))
    = \sum_{k=1}^{N}\,
    t_{ik}(u) \otimes t_{kj}(u - z),
  \end{equation}
  where $t_{ij}(u) = \delta_{ij} + \sum_{r \geq 1}
  t_{ij}^{(r)}\, u^{-r}$.
  This is coassociative
  up to the Drinfeld associator $\Phi$:
  $(\Delta_{z_1} \otimes \id) \circ \Delta_{z_2}(T(u))
  = T(u) \otimes T(u - z_2) \otimes T(u - z_1 - z_2)$
  and
  $(\id \otimes \Delta_{z_2}) \circ \Delta_{z_1}(T(u))
  = T(u) \otimes T(u - z_1) \otimes T(u - z_1 - z_2)$,
  strictly coassociative for $\mathfrak{gl}_N$
  at generic $\Psi$.
\item \textup{(RTT relation: non-trivial for $N \geq 2$.)}
  The RTT relation
  \begin{equation}\label{eq:glN-rtt}
    R_{12}(u - v)\, T_1(u)\, T_2(v)
    = T_2(v)\, T_1(u)\, R_{12}(u - v),
  \end{equation}
  where $T_1(u) = T(u) \otimes I_N$,
  $T_2(v) = I_N \otimes T(v)$, is the defining relation of
  $Y(\mathfrak{gl}_N)$.
  For $N = 1$, this is trivially satisfied
  \textup{(}scalar $R$\textup{)}. For $N \geq 2$ it is a
  genuine quadratic constraint on the generators
  $t_{ij}^{(r)}$: the commutator $[R_{12}, T_1 T_2]$
  is generically nonzero
  \textup{(}verified numerically at $N = 2, 3$\textup{)}.
\item \textup{(Quantum determinant.)}
  The quantum determinant
  \begin{equation}\label{eq:glN-qdet}
    \mathrm{qdet}\, T(u)
    = \sum_{\sigma \in S_N}
    \mathrm{sgn}(\sigma)\,
    \prod_{i=1}^{N}
    T_{i,\sigma(i)}\bigl(u - (i-1)\Psi\bigr)
  \end{equation}
  is \emph{central} in $Y(\mathfrak{gl}_N)$: it commutes
  with all $t_{ij}(v)$. In the evaluation representation,
  $\mathrm{qdet}\, T(u)$ is proportional to the identity
  on~$\CC^N$ \textup{(}verified numerically at
  $N = 2, 3$\textup{)}.
  The $\Psi$-shift $u \mapsto u - (i-1)\Psi$
  is the quantum correction to the classical determinant
  $\det T(u)$.
\item \textup{($\Ainf$ structure.)}
  The chiral $\Ainf$-operations $m_k^{\mathrm{ch}}$ on
  $\cW_N[\Psi]$ are extracted from the ordered bar complex
  $\Barord(\cW_N) = T^c(s^{-1}\overline{\cW}_N)$ via the
  functor $\alpha$ of
  Theorem~\textup{\ref{thm:chiral-qg-equiv}}, with
  input the $R$-matrix~\eqref{eq:glN-yang-r}.
\end{enumerate}
The three structures \textup{(}$R$-matrix, coproduct,
$\Ainf$\textup{)} determine each other by
Theorem~\textup{\ref{thm:chiral-qg-equiv}}.
The OPE compatibility axiom~\eqref{eq:ope-compat}
is satisfied at all spins $1, \ldots, N$ simultaneously
by the same two independent arguments as in
Theorem~\textup{\ref{thm:w-infty-chiral-qg}}:
\begin{enumerate}[label=\textup{(\alph*)}]
\item \emph{Argument~$A$ \textup{(}coderivation on the
  Koszul locus\textup{)}.}
  At generic~$\Psi$, the algebra $\cW_N[\Psi]$ lies on the
  Koszul locus
  \textup{(}Definition~\textup{\ref{def:koszul-locus}}\textup{)}.
  The bar-cobar adjunction
  $\Omega(\Barord(\cW_N[\Psi]))
  \xrightarrow{\;\sim\;} \cW_N[\Psi]$
  is an equivalence, and the coderivation property
  $d_{\Barord} \circ \Delta = \Delta \circ d_{\Barord}$
  on $\Barord(\cW_N[\Psi]) = T^c(s^{-1}\overline{\cW}_N[\Psi])$
  dualizes to OPE
  compatibility~\eqref{eq:ope-compat} at all spins
  simultaneously
  \textup{(}Theorem~\textup{\ref{thm:chiral-qg-equiv}},
  direction $\textup{(II)} \to \textup{(III)}$\textup{)}.
  Uniqueness on the Koszul locus identifies this bar-cobar
  coproduct with the Drinfeld
  coproduct~\eqref{eq:glN-drinfeld-coprod}.
\item \emph{Argument~$B$ \textup{(}JKL vertex bialgebra
  on the CoHA\textup{)}.}
  The critical CoHA of the Jordan quiver at dimension
  vector $n \cdot \delta$ ($\delta$ the minimal imaginary
  root of $\widehat{\mathfrak{gl}}_N$) satisfies the vertex
  bialgebra axiom~\eqref{eq:ope-compat} by
  Jindal--Kaubrys--Latyntsev~\textup{\cite{JKL26}},
  Theorem~C for general quivers. The
  Schiffmann--Vasserot identification
  $\cH_{\mathrm{Jor},N} \cong Y^+(\widehat{\mathfrak{gl}}_N)$
  transports this to the Yangian, giving a second proof
  at all spins simultaneously.
\end{enumerate}
\end{theorem}

\begin{proof}
The proof follows the seven-step structure of
Theorem~\ref{thm:w-infty-chiral-qg}, with the scalar
$R$-matrix replaced by Yang's matrix $R$-matrix.

\medskip
\noindent
\textbf{Step~1}
(CoHA identification; Schiffmann--Vasserot~\cite{SV13}).
The critical CoHA of the framed Jordan quiver at
rank~$N$ satisfies
$\cH_{\mathrm{Jor},N} \cong Y^+(\widehat{\mathfrak{gl}}_N)$.
The $N = 1$ case is
Schiffmann--Vasserot~\cite{SV13}; the general case
uses the Drinfeld--Sokolov hierarchy for
$\widehat{\mathfrak{gl}}_N$.

\medskip
\noindent
\textbf{Step~2}
(Yangian--$\cW_N$ identification;
Proch\'azka--Rap\v{c}\'ak~\cite{PR19}).
The quantum Miura transform gives an isomorphism of
filtered associative algebras
$Y(\widehat{\mathfrak{gl}}_N) \cong U(\cW_N[\Psi])$.
The transfer matrix~\eqref{eq:glN-transfer}
maps to the Miura
factorisation~\eqref{eq:glN-miura}.
At $N = 2$: $\cW_2 = \mathrm{Vir}_c$ with
$c = 1 - 6(\Psi - 1)^2/\Psi$ and
$\kappa(\mathrm{Vir}_c) = c/2$.
At $N = 3$: $\cW_3$ has generators at spins $2, 3$ and
$\kappa(\cW_3) = 5c/6$ (where $H_3 - 1 = 5/6$).

\medskip
\noindent
\textbf{Step~3}
($R$-matrix: Yang's rational $R$-matrix).
The $R$-matrix~\eqref{eq:glN-yang-r} in the fundamental
representation of $\mathfrak{gl}_N$ is:
\[
  R(u) = u\, I_{N^2} + \Psi\, P
  \qquad \textup{(additive convention)}.
\]
% AP126: level prefix Psi; Psi=0 -> R = u I (trivial). Verified.
The level prefix $\Psi$ is mandatory: at $\Psi = 0$,
$R(u) = u\, I$ (no interaction). The classical
$r$-matrix is $r(z) = \Psi\, P/z$, with
$r|_{\Psi = 0} = 0$.
The key properties, verified numerically at $N = 2, 3$ and
proved algebraically for all~$N$
(Molev~\cite{Molev07}):
\begin{enumerate}[label=\textup{(\roman*)}]
\item Permutation involution: $P^2 = I$.
\item Unitarity: $R(u)\, R(-u) = (\Psi^2 - u^2)\, I$.
\item Yang--Baxter equation~\eqref{eq:glN-ybe}.
\end{enumerate}
At $N = 1$: $P = 1$ on $\CC^1 \otimes \CC^1$ and
$R(u) = u + \Psi$ is scalar, so the RTT
relation~\eqref{eq:glN-rtt} is trivially satisfied.
For $N \geq 2$: $P$ is a genuine $N^2 \times N^2$
permutation matrix with eigenvalues $\pm 1$, and the
RTT is non-trivial.

\medskip
\noindent
\textbf{Step~4}
(Drinfeld coproduct at all spins).
The matrix Drinfeld
formula~\eqref{eq:glN-drinfeld-coprod} gives
$\Delta_z(T(u)) = T(u) \cdot T(u - z)$. In components,
the coefficient of $u^{-n}$ is:
\[
  \Delta_z(T^{(n)})_{ij}
  = \sum_{\substack{a,b \geq 0 \\ a+b+c = n}}
  \binom{b+c-1}{c}\, z^c\,
  \sum_{k=1}^{N}
  t_{ik}^{(a)} \otimes t_{kj}^{(b)},
\]
where $c = n - a - b \geq 0$ and
$\binom{-1}{0} = 1$ by convention. The matrix structure
sums over the auxiliary index $k$, distinguishing this from
the scalar ($N = 1$) formula~\eqref{eq:gl1-coprod-general}.

\medskip
\noindent
\textbf{Step~5}
(RTT verification).
The RTT relation~\eqref{eq:glN-rtt} is a consequence of
the Yang--Baxter equation~\eqref{eq:glN-ybe} by
evaluating the universal $R$-matrix in the
fundamental representation. In the evaluation
representation $T(u) \mapsto R(u - a)$, the RTT becomes
the YBE with shifted spectral parameters.
The RTT at
$N \geq 2$ encodes genuine quadratic relations among the
generators $t_{ij}^{(r)}$: the commutator
$[R_{12}(u-v),\, T_1(u)\, T_2(v)]$ is generically nonzero
(verified at $N = 2$: commutator norm $= 1.0$;
$N = 3$: commutator norm $= 1.0$;
$N = 1$: commutator $= 0$).

\medskip
\noindent
\textbf{Step~6}
(OPE compatibility at all spins: two arguments).
The two independent proofs of OPE compatibility extend
from $\cW_{1+\infty}$ to $\cW_N$:

\smallskip
\noindent\emph{Argument~A.}
The coderivation property
$d_{\Barord} \circ \Delta = \Delta \circ d_{\Barord}$
on the bar complex $\Barord(\cW_N[\Psi])
= T^c(s^{-1}\overline{\cW}_N[\Psi])$
is a structural property of the cofree tensor coalgebra,
valid at each truncation $N$. At generic~$\Psi$,
$\cW_N[\Psi]$ lies on the Koszul locus, and the
bar-cobar adjunction identifies the bar-cobar coproduct
with the Drinfeld coproduct by uniqueness
(Theorem~\ref{thm:chiral-qg-equiv}).

\smallskip
\noindent\emph{Argument~B.}
The JKL vertex bialgebra theorem~\cite{JKL26} applies to
the CoHA of the Jordan quiver at arbitrary rank~$N$
(their Theorem~C holds for general quivers with
potential). The vertex coproduct on
$\cH_{\mathrm{Jor},N}
\cong Y^+(\widehat{\mathfrak{gl}}_N)$
satisfies~\eqref{eq:ope-compat} and identifies with
the Drinfeld coproduct.

\medskip
\noindent
\textbf{Step~7}
(Quantum determinant and full axiom verification).
The quantum determinant~\eqref{eq:glN-qdet} is central
in $Y(\mathfrak{gl}_N)$
(Molev~\cite{Molev07}, Theorem~1.11.2),
verified numerically to be proportional to
$I_N$ in the evaluation representation at $N = 2, 3$.
The coproduct satisfies coassociativity, the counit axiom,
and OPE compatibility at all spins $1, \ldots, N$ by
step~6. Together with
the Yang $R$-matrix and the $\Ainf$ structure from the bar
complex, this gives the complete chiral quantum group
datum for $\cW_N[\Psi]$.
\end{proof}

\begin{remark}[The $N = 1 \to N \geq 2$ transition]
\label{rem:glN-transition}
The passage from $N = 1$
(Theorem~\ref{thm:w-infty-chiral-qg}) to $N \geq 2$
(Theorem~\ref{thm:glN-chiral-qg}) has the following
structural features.
\begin{enumerate}[label=\textup{(\arabic*)}]
\item The $R$-matrix transitions from scalar
  ($R(u) = u + \Psi$, trivialising the RTT) to
  matrix-valued ($R(u) = u\, I + \Psi\, P$, making
  the RTT a genuine constraint). This is the
  non-abelian threshold.
\item The transfer matrix transitions from the scalar
  series $T(u) = 1 + \sum \psi_n u^{-n}$ to the
  $N \times N$ matrix
  $T(u) = I_N + \sum T^{(r)} u^{-r}$.
  The Drinfeld coproduct
  $\Delta_z(T(u)) = T(u) \cdot T(u - z)$ is matrix
  multiplication in $\CC^N$, not scalar multiplication.
\item The quantum determinant $\mathrm{qdet}\, T(u)$ is a
  genuinely new central element for $N \geq 2$ (at
  $N = 1$ it reduces to $T(u)$ itself).
\item Both arguments for OPE compatibility (coderivation on
  the Koszul locus, JKL on the CoHA) apply uniformly for
  all~$N$: argument~A uses only the structural properties
  of the cofree tensor coalgebra
  $T^c(s^{-1}\overline{\cW}_N)$ and the Koszul-locus
  equivalence; argument~B uses JKL for the Jordan quiver
  at rank~$N$.
\end{enumerate}
\end{remark}


% ----------------------------------------------------------------
\subsection{Structural consequences}
\label{subsec:structural-consequences}

\begin{corollary}[The ordered bar complex encodes all three
structures]
\label{cor:bar-encodes-all}
The ordered bar complex $\Barord(\cA) = T^c(s^{-1}\bar{\cA})$,
equipped with its bar differential $d_{\Barord}$ and
deconcatenation coproduct $\Delta$, is the universal
object from which all three structures in
Theorem~\textup{\ref{thm:chiral-qg-equiv}} are recovered:
\begin{enumerate}[label=\textup{(\roman*)}]
\item The bar differential $d_{\Barord}$ encodes the
  $\Ainf$ structure maps $m_k^{\mathrm{ch}}$
  \textup{(}by the standard bar-cobar correspondence:
  coderivations of $T^c(s^{-1}\bar{\cA})$ biject with
  elements of $\prod_k \End^{\mathrm{ch}}_\cA(k)$\textup{)}.
\item The deconcatenation coproduct $\Delta$ induces,
  via cobar dualization, the chiral coproduct
  $\Delta^{\mathrm{ch}}$.
\item The vertex $R$-matrix $S(z)$ is recovered as the
  monodromy of the KZ connection, which is itself the
  degree-$2$ projection of the Maurer--Cartan element
  $\Theta_\cA$ in $\gEone$.
\end{enumerate}
This is the structural content of the $\Eone$ primacy
principle: the ordered bar complex with deconcatenation
is the primitive object; the $R$-matrix, the chiral
coproduct, and the $\Ainf$ operations are projections of
this single datum.
\end{corollary}

\begin{remark}[Relation to factorization quantum groups]
\label{rem:factorization-qg}
When the hexagon holds strictly, the chiral quantum group
$(\cA, S(z), \Delta^{\mathrm{ch}})$ is a
\emph{factorization quantum group}
(Latyntsev~\cite{Latyntsev23}): its ordered bar complex
defines a factorization coalgebra on the curve. When the
hexagon holds only up to $\Phi$, one obtains a
\emph{quasi-factorization quantum group} with regularized
holonomy data $(R, \Phi)$.
\end{remark}


